\section{Zone endommagée autour d'un tunnel : reproduction de Wang et al. 2024}
\label{sec-tunnel}

rockim reproduit les grandeurs globales du tunnel de Wang et al. : rayon de la zone endommagée (EDZ), prédominance du cisaillement, fissures conjuguées, effet du rapport $\lambda$. Le rayon dépend toutefois de l'estimateur : l'écart à l'article va de $-4$ à $-21~\%$ au maximum et de $-19$ à $-31~\%$ au quantile $95~\%$ entre 4 et 7~MPa. Il dépend aussi d'un amortissement choisi par rockim, qui change le nombre de joints rompus d'un facteur 4. La texture de la rupture n'est pas reproduite : rockim granule un anneau là où MultiFracS découpe des blocs. Tous les calculs, référence de production comprise, emploient la décharge \cle{origin} sans cliquet, non conservative ; en production, le travail du contact valait déjà 416~kJ/m pour 545~kJ/m de fissuration, et le résidu du bilan d'énergie dépasse le critère de $1~\%$. Le rayon d'EDZ comparé à Wang est donc lui aussi conditionnel. Un diagnostic sur maillage réduit attribue l'injection par le contact à une erreur d'intégration du potentiel de contact, qui suit sa raideur normale et le pas de temps, et non à cette décharge (\cref{sec-tunnel-diag}). Le faciès dépend lui aussi de cette raideur ; une raideur dix fois plus faible reste à vérifier avec frottement. Les variantes qui découpent des blocs, où le contact injecte plus que l'énergie cohésive, sont les plus suspectes.

Wang et al.~\cite{wang2024tunnel} simulent avec MultiFracS l'excavation du tunnel routier de Hutou Beishan, creusé dans une roche tendre sous forte contrainte in situ. Leur cas de base, sous $\sigma_0 = 5$~MPa hydrostatique, donne une EDZ quasi circulaire, des fissures conjuguées en X où domine le cisaillement, et trois modes de ruine : compression-cisaillement aux reins, extrusion en voûte et soulèvement du radier. Il n'existe pas de solution analytique pour ce faciès ; seuls l'état élastique initial (Kirsch) et un rayon plastique de Mohr-Coulomb en ont une. Le banc est donc une reproduction de code à code, comme celui de Yan (\cref{sec-yan}).

Le banc vérifie deux capacités ajoutées pour l'occasion : la précontrainte in situ (\cle{insituSh}, \cle{insituSv}) et l'excavation par relâchement de la traction de paroi (\cle{excavRelease}). Il sollicite la chaîne cohésive complète, loi de Yan et insertion adaptative, sur un problème quasi statique de grande taille, et mesure la sensibilité du faciès au schéma d'insertion, au désordre des résistances, au facteur de pointe (abaissement du seuil d'insertion devant une fissure, \cref{eq-tun-insertion}) et à l'amortissement. Les chiffres archivés datent d'août 2026 ; une partie a été rejouée le 6 octobre 2026 avec le binaire courant, sur un maillage réduit (\cref{sec-tunnel-rejeu}).

\subsection{Mise en données}

\subsubsection{Géométrie, maillage et chargement}

Le domaine est un carré de $100 \times 100$~m en déformation plane, grandeurs par mètre d'épaisseur. La section en fer à cheval, centrée en $(50, 50)$~m, a une portée de $11{,}10$~m et une hauteur de $8{,}85$~m (voûte de rayon $5{,}55$~m, congés de $1{,}30$~m, radier de $22{,}52$~m) ; son périmètre lu par le solveur vaut $33{,}21$~m. Le maillage de production est un Delaunay non structuré (Gmsh, algorithme 5), de taille $h = 0{,}22$~m jusqu'à 18~m de la paroi et 2~m au loin : $106\,298$ triangles et $159\,269$ arêtes cohésives potentielles. Wang et al. emploient $83\,458$ triangles.

L'état initial est hydrostatique, $\sigma_h = \sigma_v = \sigma_0 = 5$~MPa ($\lambda = \sigma_h/\sigma_v = 1$), sans pesanteur. Deux balayages s'y ajoutent : $\sigma_0$ de 3 à 7~MPa à $\lambda = 1$, et $\lambda$ de $0{,}5$ à $1{,}5$ à $\sigma_h = 5$~MPa. Les flancs sont sur rouleaux ; le haut et le bas sont tenus par des mors immobiles, encastrés (\cref{sec-tunnel-controles}). Le calcul s'équilibre d'abord pendant $0{,}02$~s sous la traction de paroi complète, puis la traction décroît en $0{,}08$~s, soit environ $1{,}7$ transit d'onde sur 100~m. La durée de référence est de $0{,}25$~s, portée à $0{,}7$ ou $1{,}0$~s pour les variantes.

\subsubsection{Excavation : écart de procédure avec l'article}

Wang et al. donnent au noyau une résistance artificielle pendant l'équilibrage, puis réduisent progressivement son module et son amortissement, méthode de réduction du module du noyau. rockim maille le massif avec sa cavité, applique $\sigma_0$ comme contrainte initiale de chaque élément et rétablit sur chaque face de paroi la traction qu'exerçait la roche excavée, avant de la faire décroître à zéro, méthode de convergence-confinement :
\begin{gather}
\boldsymbol\sigma_e = \boldsymbol\sigma_0 + \mathbf D:\boldsymbol\varepsilon(\mathbf u), \qquad \boldsymbol\sigma_0 = -\operatorname{diag}(\sigma_h, \sigma_v), \label{eq-tun-insitu}\\
\mathbf t_\text{paroi}(t) = r(t)\,\boldsymbol\sigma_0\cdot\mathbf n, \qquad
r(t) = \begin{cases} 1 & t \leq t_s, \\ \tfrac12\left[1 + \cos\dfrac{\pi(t - t_s)}{t_r}\right] & t_s < t < t_s + t_r, \\ 0 & t \geq t_s + t_r, \end{cases} \label{eq-tun-relache}
\end{gather}
avec $t_s = 0{,}02$~s (\cle{excavStart}) et $t_r = 0{,}08$~s (\cle{excavRamp}). Les deux procédures ont le même état initial et le même état final ; celle de rockim ne fait varier aucun module, et la traction face par face reste exacte quand $\lambda \neq 1$.

L'insertion adaptative est ici nécessaire. Tant que le critère n'est pas atteint, les sommets des arêtes restent des nœuds partagés et la précontrainte des éléments décrit exactement l'état initial. En insertion intrinsèque, chaque joint devrait d'abord se fermer de $\sigma_0/p_j$ sous la précontrainte ; le calcul intrinsèque de contrôle décale donc le relâchement à $t_s = 0{,}15$~s pour laisser cette fermeture s'amortir.

\subsubsection{Paramètres}

Les paramètres du matériau et des joints sont ceux du tableau~1 de Wang et al., sans ajustement (\cref{tab-tun-param}). L'amortissement, que l'article ne chiffre pas, est un choix de rockim, et il pilote le nombre de joints rompus d'un facteur 4 (\cref{sec-tunnel-duree}). La longueur cohésive en mode~I vaut $\ell_{cz} = E\,G_I/f_t^2 = 0{,}556$~m et celle du mode~II $E\,G_{II}/c^2 = 1{,}56$~m. La règle d'objectivité $h < \ell_{cz}/2 = 0{,}278$~m est satisfaite par $h = 0{,}22$~m. L'enveloppe de Mohr-Coulomb des joints donne une résistance en compression $2c\cos\varphi/(1 - \sin\varphi) = 3{,}0$~MPa ; une éprouvette numérique à la taille de maille de la zone fine donne $3{,}94$~MPa. Le rapport $R_c/\sigma_0 = 5$ que cite l'article est un indice de classification de la roche intacte, pas la résistance du modèle des auteurs.

\begin{table}[htbp]
\centering
\caption{Paramètres du banc du tunnel : tableau~1 de Wang et al. et réglage de rockim.}
\label{tab-tun-param}
\small
\begin{tabular}{@{}p{0.25\linewidth}p{0.33\linewidth}p{0.36\linewidth}@{}}
\toprule
paramètre & Wang et al. & rockim \\
\midrule
$\rho$, $E$, $\nu$ & 2\,650~kg/m$^3$, 10~GPa, $0{,}26$ & identiques \\
$f_t$, $c$, $\varphi$ (joints) & $0{,}6$~MPa, $0{,}8$~MPa, $34^\circ$ & identiques \\
$G_I$, $G_{II}$ & 20 et 100~J/m$^2$ & identiques (\cle{gfShearFactor = 5}) \\
pénalité des joints & 1\,000~GPa, soit $100\,E$ & $4\,E/h$ à l'insertion \\
adoucissement & courbe de Munjiza & \cle{jointSoftening = yan} \\
décharge en cisaillement & non précisée & \cle{jointShearUnload = origin}, sans cliquet (\cle{jointSecantRatchet} désactivé) \\
éléments & élastiques & élastiques (plafond déviatorique neutralisé) \\
contact & potentiel & potentiel de Munjiza \\
pas de temps & $\approx 8\times10^{-7}$~s, $300\,000$ pas & $\approx 3{,}05\times10^{-6}$~s (\cle{dtFactor = 0.2}) \\
amortissement & visqueux, non chiffré & local de Cundall $0{,}15$, $\xi_j = 0{,}01$ (\cle{jointXi}) \\
\bottomrule
\end{tabular}
\end{table}

\subsection{Équations du banc}

\subsubsection{Solution de Kirsch}

Pour un trou circulaire de rayon $a$ dans un milieu infini sous les compressions lointaines $\sigma_h$ et $\sigma_v$ (comptées positives), avec $\theta$ mesuré depuis l'horizontale :
\begin{gather}
\sigma_{rr} = \frac{\sigma_h + \sigma_v}{2}\left(1 - \frac{a^2}{r^2}\right) + \frac{\sigma_h - \sigma_v}{2}\left(1 - \frac{4a^2}{r^2} + \frac{3a^4}{r^4}\right)\cos 2\theta, \\
\sigma_{\theta\theta} = \frac{\sigma_h + \sigma_v}{2}\left(1 + \frac{a^2}{r^2}\right) - \frac{\sigma_h - \sigma_v}{2}\left(1 + \frac{3a^4}{r^4}\right)\cos 2\theta, \\
\sigma_{r\theta} = -\frac{\sigma_h - \sigma_v}{2}\left(1 + \frac{2a^2}{r^2} - \frac{3a^4}{r^4}\right)\sin 2\theta .
\end{gather}
En paroi, $\sigma_{\theta\theta}(a, \theta) = (\sigma_h + \sigma_v) - 2(\sigma_h - \sigma_v)\cos 2\theta$ : 10~MPa partout pour $\lambda = 1$ et $\sigma_0 = 5$~MPa ; 5~MPa en couronne et 25~MPa aux reins pour $\sigma_h = 5$ et $\sigma_v = 10$~MPa. La tolérance retenue est de $5~\%$, la plaque étant finie ($a/b = 1/10$, $b$ demi-côté du domaine).

\subsubsection{Critère d'insertion, facteur de pointe et endommagement}

Une arête liée est insérée quand la traction moyenne des deux éléments voisins, projetée sur l'arête en configuration déformée, franchit l'enveloppe de Yan et al. 2023 :
\begin{equation}
\sigma_n \geq \frac{f_t}{k} \quad \text{ou} \quad |\tau| \geq \frac{c + \tan\varphi\,\langle -\sigma_n\rangle}{k},
\label{eq-tun-insertion}
\end{equation}
où $k = k_\text{tip}$ si l'arête touche un joint d'endommagement $D \geq D_\text{tip}$, et $k = 1$ sinon. Le facteur de pointe $k_\text{tip}$ (\cle{insertionTipFactor}, avec $D_\text{tip} = 0{,}5$ par \cle{insertionTipDamage}) est une option de rockim : il abaisse le seuil devant une fissure qui s'ouvre et favorise ainsi la propagation plutôt que la nucléation. Il vaut 1 dans le calcul de référence. L'endommagement du joint suit la forme de Munjiza reprise par Wang et al., avec $o$ et $s$ les ouvertures normale et tangentielle, $o_p$ et $s_p$ leurs valeurs au pic, $o_r$ et $s_r$ à la rupture :
\begin{equation}
D = \min\left\{1,\ \sqrt{\left\langle\frac{o - o_p}{o_r - o_p}\right\rangle^2 + \left\langle\frac{s - s_p}{s_r - s_p}\right\rangle^2}\right\},
\label{eq-tun-D}
\end{equation}
qui redonne les branches en mode pur de l'article quand l'un des deux termes est nul ; c'est la forme elliptique de~\eqref{eq-D} écrite à partir du pic.

\subsubsection{Mesure de l'EDZ}

Wang et al. définissent l'EDZ comme la région d'expansion des fissures et la lisent sur leurs figures. rockim la mesure sur l'ensemble $\mathcal B$ des joints rompus ($D \geq 1$), de milieux $\mathbf x_j$ et de distances $r_j = \|\mathbf x_j - \mathbf x_c\|$ au centre du tunnel :
\begin{equation}
R_\text{EDZ}^\text{max} = \max_{j \in \mathcal B} r_j, \qquad
R_\text{EDZ}^{95} = Q_{0,95}\{r_j\}, \qquad
a_h = Q_{0,95}\{|x_j - x_c|\}, \quad a_v = Q_{0,95}\{|y_j - y_c|\},
\label{eq-tun-edz}
\end{equation}
où $Q_{0,95}$ est le quantile à $95~\%$. Le maximum est l'estimateur le plus proche d'une lecture visuelle, mais une seule fissure isolée le fixe ; le quantile est plus robuste. Un joint rompu est classé mixte si $\min(r_n, r_s)/\max(r_n, r_s) \geq 0{,}5$, où $r_n$ et $r_s$ sont les parts normale et tangentielle de l'endommagement, sinon en traction ou en cisaillement selon le rapport dominant. La convergence de paroi est le déplacement moyen des nœuds de paroi, insensible aux blocs détachés, contrairement au maximum du champ.

\subsubsection{Rayon plastique de Mohr-Coulomb}

Pour un trou circulaire sous pression hydrostatique, paroi non soutenue, dans un milieu élastique parfaitement plastique de Mohr-Coulomb, avec $K_p = (1 + \sin\varphi)/(1 - \sin\varphi)$ :
\begin{equation}
R_p = a\left[\frac{2}{K_p + 1}\,\frac{\sigma_0 + c\cot\varphi}{c\cot\varphi}\right]^{1/(K_p - 1)} .
\label{eq-tun-rp}
\end{equation}
Avec $c = 0{,}8$~MPa, $\varphi = 34^\circ$ ($K_p = 3{,}54$) et le rayon équivalent $a = 33{,}21/(2\pi) = 5{,}29$~m, $R_p$ vaut $6{,}3$, $6{,}9$, $7{,}3$, $7{,}8$ et $8{,}2$~m pour $\sigma_0 = 3$ à 7~MPa. Les EDZ de rockim et de Wang sont deux à trois fois plus grandes : le massif FDEM est fragile, il perd toute cohésion après rupture, et l'EDZ compte tout joint rompu, pas la seule frontière plastique. $R_p$ n'est donc qu'une borne basse d'ordre de grandeur.

\subsection{Contrôles analytiques}
\label{sec-tunnel-controles}

Avant tout calcul de rupture, les deux capacités ajoutées ont été vérifiées sur des cas à réponse connue (\cref{tab-tun-controles}). Le calcul à charge nulle, où l'excavation n'est jamais déclenchée, doit rester immobile : il a d'abord donné $5\,185$~J/m d'énergie cinétique parasite, venue des quatre nœuds de coin laissés libres en $x$ sous $\sigma_{xx}$, chacun portant environ $7{,}5$~MN. Les mors encastrés ont corrigé ce défaut ; l'énergie cinétique résiduelle est de $2\times10^{-17}$~J/m. Le contrôle de Kirsch montre que ce choix de bord ne modifie pas l'état en paroi. En régime quasi statique, le résidu du bilan d'énergie n'est pas lisible : la correction d'intégration aux nœuds à vitesse imposée vaut $13{,}6$~J/m par pas. Le jugement porte donc sur l'énergie cinétique et sur une jauge de contrainte.

\begin{table}[htbp]
\centering
\caption{Contrôles analytiques de la précontrainte et de l'excavation (août 2026).}
\label{tab-tun-controles}
\small
\begin{tabular}{@{}p{0.30\linewidth}p{0.27\linewidth}p{0.36\linewidth}@{}}
\toprule
contrôle & attendu & mesuré \\
\midrule
suite rapide après ajout des deux options & résultats identiques au bit près & 15 cas sur 15 identiques \\
charge nulle, excavation non déclenchée & aucun joint inséré, jauge à 5~MPa & aucune insertion, jauge $5{,}000$~MPa, $E_c = 2\times10^{-17}$~J/m \\
Kirsch, $\lambda = 1$, $a = 5$~m, joints incassables & $\sigma_{\theta\theta} = 10$~MPa en paroi & $9{,}84$~MPa extrapolé, écart $1{,}7~\%$ \\
Kirsch, $\lambda = 0{,}5$ & 5~MPa en couronne, 25~MPa aux reins & $5{,}51$ et $24{,}05$~MPa pour $5{,}49$ et $24{,}57$~MPa au même rayon, écart $2{,}1~\%$ \\
résistance en compression à la maille fine & au-dessus de $3{,}0$~MPa & $3{,}94$~MPa \\
orientation des fissures conjuguées & $\pm(45^\circ - \varphi/2) = \pm 28^\circ$ & surreprésentation $\times 2{,}00$ à $-25^\circ$ et $\times 1{,}91$ à $+25^\circ$ \\
\bottomrule
\end{tabular}
\end{table}

\subsection{Cas de référence}

Le tableau~\ref{tab-tun-ref} compare le cas de référence à l'article. Le rayon d'EDZ maximal est de $17{,}19$~m contre 19~m lus chez Wang ($-9{,}5~\%$), et le rayon au quantile $95~\%$ de $14{,}49$~m ($-24~\%$). Le maximum est l'estimateur le plus favorable à l'accord, et les valeurs de Wang, lues sur des figures, n'ont pas d'incertitude annoncée. L'EDZ est quasi circulaire, de demi-axes $12{,}88$ et $12{,}52$~m au quantile $95~\%$ (\cref{fig-tun-wang11,fig-tun-fissures}). Le déplacement maximal, $0{,}389$~m, se trouve sur un bloc détaché ; la convergence moyenne de paroi n'est que de $0{,}069$~m. Les sources de Wang divergent : on lit $0{,}347$~m sur la figure~14a, le texte annonce $0{,}373$~m ; l'écart est de $+12~\%$ ou de $+4~\%$ selon la valeur retenue. Le cisaillement domine, comme chez Wang, mais la traction passe devant le mode mixte, inversant l'ordre publié ; ce classement dépend du seuil arbitraire $0{,}5$. Les fissures conjuguées se forment à $\pm 25^\circ$, près de l'angle de Mohr-Coulomb.

\begin{table}[htbp]
\centering
\caption{Cas de référence ($\sigma_0 = 5$~MPa, $\lambda = 1$, $t = 0{,}25$~s) : rockim comparé à Wang et al. 2024. Les valeurs de Wang sont lues sur les figures.}
\label{tab-tun-ref}
\small
\begin{tabular}{@{}p{0.2\linewidth}p{0.32\linewidth}p{0.42\linewidth}@{}}
\toprule
grandeur & Wang et al. & rockim \\
\midrule
rayon d'EDZ & 19~m (fig.~12f) & $17{,}19$~m (max), $14{,}49$~m (quantile $95~\%$) \\
forme & quasi circulaire & demi-axes $12{,}88$ et $12{,}52$~m \\
déplacement maximal & $0{,}347$~m (fig.~14a), $0{,}373$~m (texte) & $0{,}389$~m sur un bloc ; paroi $0{,}069$~m en moyenne \\
joints rompus & non publié, environ $10\,000$ sur la fig.~8 & $14\,935$ \\
modes & cisaillement, puis mixte, puis traction & cisaillement $7\,747$, traction $4\,817$, mixte $2\,371$ \\
longueur fissurée & non publiée & $3\,115$~m \\
faciès & X conjugué & X conjugué à $\pm 25^\circ$ \\
calcul & GPU, $300\,000$ pas & 44~min sur processeur multifil \\
\bottomrule
\end{tabular}
\end{table}

\begin{figure}[htbp]
\centering
\includegraphics[width=\linewidth]{tunnel/fig_wang11}
\caption{Reproduction de la figure~11 de Wang et al. : déplacement en configuration déformée et fissures à $t = 0{,}25$~s ; zooms : serrage aux reins, compression-cisaillement, soulèvement du radier. Maximum de paroi hors débris $0{,}148$~m, $0{,}389$~m sur un bloc détaché.}
\label{fig-tun-wang11}
\end{figure}

\begin{figure}[htbp]
\centering
\includegraphics[width=\linewidth]{tunnel/ref_fissures}
\caption{Réseau de fissures du cas de référence, vue d'ensemble et zoom sur la paroi. La figure classe en deux modes, selon la part dominante (traction en vert, 6\,289 ; cisaillement en rouge, 8\,645), et non en trois comme le texte. Anneau quasi circulaire, granulé près de la paroi.}
\label{fig-tun-fissures}
\end{figure}

\subsection{Balayages de la contrainte in situ et du rapport \texorpdfstring{$\lambda$}{lambda}}

Le balayage de $\sigma_0$ suit Wang et al. de 3 à 5~MPa, puis sature (\cref{tab-tun-s0,fig-tun-sweep}). Les écarts sur le rayon maximal vont de $+15~\%$ à 3~MPa à $-21~\%$ à 6~MPa ; au quantile $95~\%$, de $-1~\%$ à $-31~\%$. Le nombre de joints rompus, la longueur fissurée et la convergence de paroi croissent de façon monotone. Le cisaillement prend le dessus et croît le plus vite (de $1\,358$ à $14\,250$ joints) tandis que la traction sature, comme sur la figure~13a de l'article. La saturation de l'EDZ au-delà de 5~MPa, absente chez Wang, n'est pas expliquée et reste un point ouvert.

\begin{table}[htbp]
\centering
\caption{Balayage de la contrainte in situ ($\lambda = 1$, $t = 0{,}25$~s). Rayon de Wang et al. lu sur leur figure~12f ; écarts au maximum et au quantile $95~\%$.}
\label{tab-tun-s0}
\small
\setlength{\tabcolsep}{4pt}
\begin{tabular}{@{}ccccccccc@{}}
\toprule
$\sigma_0$ & rompus & traction / mixte / cisaill. & $R^\text{max}$ & $R^{95}$ & Wang & écart & écart & paroi \\
(MPa) & & & (m) & (m) & (m) & $R^\text{max}$ & $R^{95}$ & (m) \\
\midrule
3 & $3\,356$ & $1\,444 / 554 / 1\,358$ & $12{,}77$ & $11{,}02$ & $11{,}1$ & $+15~\%$ & $-1~\%$ & $0{,}018$ \\
4 & $9\,123$ & $3\,288 / 1\,470 / 4\,365$ & $15{,}42$ & $12{,}97$ & $16{,}0$ & $-4~\%$ & $-19~\%$ & $0{,}044$ \\
5 & $14\,935$ & $4\,817 / 2\,371 / 7\,747$ & $17{,}19$ & $14{,}49$ & $19{,}0$ & $-9{,}5~\%$ & $-24~\%$ & $0{,}069$ \\
6 & $20\,028$ & $6\,110 / 3\,006 / 10\,912$ & $17{,}40$ & $15{,}29$ & $22{,}0$ & $-21~\%$ & $-31~\%$ & $0{,}092$ \\
7 & $24\,535$ & $6\,808 / 3\,477 / 14\,250$ & $18{,}59$ & $15{,}67$ & $22{,}8$ & $-18~\%$ & $-31~\%$ & $0{,}119$ \\
\bottomrule
\end{tabular}
\end{table}

\begin{figure}[htbp]
\centering
\includegraphics[width=\linewidth]{tunnel/sweep_insitu}
\caption{Balayage de la contrainte in situ ($\lambda = 1$) : (a) rayon d'EDZ, maximum et quantile $95~\%$, et valeurs lues chez Wang et al. ; (b) convergence : maximum du champ, quantile à $90~\%$ et moyenne de paroi ; (c) modes ; (d) longueur fissurée.}
\label{fig-tun-sweep}
\end{figure}

À $\sigma_h = 5$~MPa et $\sigma_v = 5/\lambda$, le nombre de joints rompus décroît de $24\,297$ à $\lambda = 0{,}5$ à $9\,069$ à $\lambda = 1{,}5$, soit un rapport de $2{,}7$ ; Wang et al. passent d'environ $14\,000$ à environ $5\,000$, même rapport (\cref{tab-tun-lambda,fig-tun-lambda}). L'EDZ bascule d'une ellipse horizontale à une ellipse verticale en traversant $\lambda = 1$, comme dans l'article. La comparaison des axes reste qualitative : la figure trace $2a_h$ et $2a_v$, alors que Wang publie environ 18~m d'axe à $\lambda = 1$ pour un rayon de 19~m, ce qui laisse ambigu le sens de l'axe publié (longueur totale ou demi-axe).

\begin{table}[htbp]
\centering
\caption{Balayage du rapport de contraintes ($\sigma_h = 5$~MPa, $t = 0{,}25$~s).}
\label{tab-tun-lambda}
\small
\begin{tabular}{@{}cccccc@{}}
\toprule
$\lambda$ & rompus & $a_h$ (m) & $a_v$ (m) & $R^\text{max}$ (m) & $u^\text{max}$ (m) \\
\midrule
$0{,}5$ & $24\,297$ & $17{,}66$ & $9{,}42$ & $21{,}67$ & $0{,}366$ \\
$0{,}75$ & $19\,069$ & $15{,}45$ & $11{,}07$ & $19{,}08$ & $0{,}400$ \\
$1{,}0$ & $14\,935$ & $12{,}88$ & $12{,}52$ & $17{,}19$ & $0{,}389$ \\
$1{,}25$ & $11\,785$ & $10{,}32$ & $13{,}84$ & $17{,}83$ & $0{,}370$ \\
$1{,}5$ & $9\,069$ & $7{,}63$ & $14{,}65$ & $17{,}19$ & $0{,}169$ \\
\bottomrule
\end{tabular}
\end{table}

\begin{figure}[htbp]
\centering
\includegraphics[width=0.9\linewidth]{tunnel/fig_lambda}
\caption{Balayage du rapport de contraintes $\lambda$ : (a) axes de l'EDZ ($2a_h$, $2a_v$) et valeurs lues sur la figure~15f de Wang et al. ; (b) nombre de fissures de rockim, qui décroît d'un facteur $2{,}7$, comme chez Wang et al. (non tracé).}
\label{fig-tun-lambda}
\end{figure}

\subsection{Texture de la rupture : schéma d'insertion, désordre et facteur de pointe}

Les grandeurs globales sont reproduites ; la texture ne l'est pas. rockim produit un anneau granulé : trois blocs sur quatre sont formés d'un seul triangle, de côté moyen $1{,}3\,h$. MultiFracS découpe des blocs.

\subsubsection{Insertion intrinsèque et insertion adaptative}

Sur le même maillage, avec les mêmes lois et au même temps depuis le relâchement, l'insertion intrinsèque rompt $72~\%$ de joints en plus et découpe deux fois plus de blocs (\cref{tab-tun-intr,fig-tun-intr}). La différence porte surtout sur les gros blocs libres : 37~m$^2$ au plus en adaptatif, $318$, $120$, $114$, $104$ et 68~m$^2$ en intrinsèque. La part de propagation, c'est-à-dire la part des ruptures qui prolongent un joint déjà rompu, passe de $43{,}7~\%$ à $58{,}9~\%$. Le schéma d'insertion ne fixe donc pas la finesse de la granulation, que l'intrinsèque produit aussi près de la paroi : il fixe la capacité à isoler de gros blocs. Les $159\,269$ interfaces préexistantes de l'intrinsèque offrent des chemins continus qui ferment des contours. La pénalité finie de l'intrinsèque coûte aussi : pas de temps divisé par $1{,}9$ et module effectif réduit à $0{,}95\,E$.

L'écart ne vient pas de l'orientation des joints insérés, dont la cohérence directionnelle reste au niveau du hasard dans les deux schémas. Il vient de l'absence de chemins préexistants. La souplesse artificielle des interfaces intrinsèques régularise en partie la localisation en blocs, effet connu des modèles cohésifs intrinsèques~\cite{falk2001cohesive}. Cette lecture est conditionnelle à l'injection d'énergie par le contact, maximale en intrinsèque (\cref{sec-tunnel-rejeu}).

\begin{table}[htbp]
\centering
\caption{Insertion adaptative et insertion intrinsèque, même maillage et même temps depuis le relâchement.}
\label{tab-tun-intr}
\small
\begin{tabular}{@{}lll@{}}
\toprule
grandeur & adaptative & intrinsèque \\
\midrule
joints rompus & $21\,259$ & $36\,527$ ($+72~\%$) \\
blocs & $8\,079$ & $16\,430$ \\
blocs d'un seul élément & $75{,}7~\%$ & $84{,}9~\%$ \\
plus gros blocs libres & 37~m$^2$ & 318, 120, 114, 104, 68~m$^2$ \\
part de propagation & $43{,}7~\%$ & $58{,}9~\%$ \\
convergence de paroi à $0{,}68$~s, moyenne & $0{,}124$~m & $0{,}173$~m \\
convergence de paroi à $0{,}68$~s, max hors débris & $0{,}459$~m & $0{,}545$~m \\
\bottomrule
\end{tabular}
\end{table}

\begin{figure}[htbp]
\centering
\includegraphics[width=\linewidth]{tunnel/fig_wang11_intr}
\caption{Calcul intrinsèque de contrôle, même maillage et mêmes lois que la \cref{fig-tun-wang11} : macrofractures radiales continues jusqu'à 20 à 25~m et grands blocs isolés, absents en insertion adaptative. Le déplacement indiqué en titre est le maximum de paroi hors débris.}
\label{fig-tun-intr}
\end{figure}

\subsubsection{Désordre des résistances}

Deux formes de désordre ont été essayées pour supprimer la granulation. La première tire la résistance de chaque joint selon une loi de Weibull de module $m = 6$, $P_f(\sigma) = 1 - \exp[-(\sigma/\sigma_0)^m]$, indépendamment d'un joint à l'autre. La seconde emploie un champ gaussien corrélé de longueur $\theta = 1$ ou 4~m, indépendant du maillage. À $t = 0{,}41$~s, le tirage de Weibull change le nombre de joints rompus de $-5~\%$, le champ corrélé de $+7~\%$ ($\theta = 1$~m) et de $-5~\%$ ($\theta = 4$~m, à $t = 0{,}29$~s) ; la part de blocs d'un seul élément reste entre $75$ et $81~\%$. Aucune des deux ne supprime la granulation. Le champ corrélé brise en revanche la symétrie de l'EDZ : en fin de calcul ($\theta = 4$~m, $t = 0{,}7$~s), les demi-axes valent $19{,}39$ et $14{,}36$~m.

\subsubsection{Facteur de pointe}

Le facteur de pointe agit directement sur la propagation (\cref{tab-tun-tip,fig-tun-blocs,fig-tun-tip16}). Avec $k_\text{tip} = 1{,}6$, la part de propagation monte à $57{,}4~\%$, à $1{,}5$~point de l'intrinsèque, et des blocs polygonaux de 100 à 500~m$^2$ apparaissent. Le prix est un nombre de joints rompus multiplié par $1{,}85$ à temps égal et une EDZ plus étendue ($R^{95} = 24{,}4$~m à $0{,}7$~s). Au-delà, $k_\text{tip} = 2{,}0$ n'augmente plus la propagation mais accroît encore la fragmentation. Les gros blocs s'agrandissent, mais la granulation de l'anneau augmente aussi : la part de blocs d'un seul élément passe de $75{,}7$ à $88{,}3~\%$. Avec $k_\text{tip} = 1{,}6$, les trois modes de ruine décrits par Wang et al. apparaissent en rendu éléments (\cref{fig-tun-tip16}). Le facteur de pointe rapproche la texture de celle de MultiFracS sans la reproduire ; il n'est pas activé dans les autres chapitres.

\begin{table}[htbp]
\centering
\caption{Effet du facteur de pointe $k_\text{tip}$. Part de propagation sur l'intervalle $0{,}08$ à $0{,}70$~s ; décomptes de blocs à $t = 0{,}41$~s dans un rayon de 25~m ; $k_\text{tip} = 1$ est le calcul de référence ; n.m. : non mesuré.}
\label{tab-tun-tip}
\small
\begin{tabular}{@{}lccccc@{}}
\toprule
 & $k_\text{tip} = 1$ & $1{,}3$ & $1{,}6$ & $2{,}0$ & intrinsèque \\
\midrule
part de propagation & $43{,}7~\%$ & $52{,}3~\%$ & $57{,}4~\%$ & $55{,}3~\%$ & $58{,}9~\%$ \\
joints rompus ($0{,}41$~s) & $21\,378$ & $33\,417$ & $39\,500$ & $42\,049$ & n.m. \\
blocs ($0{,}41$~s) & $8\,079$ & $16\,003$ & $19\,484$ & $20\,902$ & n.m. \\
blocs d'un élément & $75{,}7~\%$ & $85{,}7~\%$ & $88{,}3~\%$ & $88{,}2~\%$ & n.m. \\
plus gros bloc détaché (m$^2$) & $34{,}8$ & $65{,}1$ & $119{,}8$ & $190{,}1$ & n.m. \\
$R^{95}$ / $R^\text{max}$ à $0{,}7$~s (m) & n.m. & $23{,}35 / 31{,}64$ & $24{,}40 / 37{,}15$ & $24{,}33 / 49{,}20$ & n.m. \\
paroi, moyenne à $0{,}7$~s (m) & $0{,}124$ & $0{,}171$ & $0{,}209$ & $0{,}265$ & n.m. \\
\bottomrule
\end{tabular}
\end{table}

\begin{figure}[htbp]
\centering
\includegraphics[width=\linewidth]{tunnel/fig_blocs_compare}
\caption{Taille des blocs intacts (aire en échelle logarithmique) pour l'insertion adaptative, l'adaptative avec $k_\text{tip} = 1{,}6$ et l'intrinsèque : l'adaptative granule un anneau, les deux autres découpent des blocs polygonaux de quelques dizaines à quelques centaines de m$^2$.}
\label{fig-tun-blocs}
\end{figure}

\begin{figure}[htbp]
\centering
\includegraphics[width=\linewidth]{tunnel/fig_blocs_tip16}
\caption{Modes de ruine en rendu éléments, $k_\text{tip} = 1{,}6$, $t = 0{,}7$~s (les vides sont en blanc) : extrusion en voûte, compression-cisaillement au rein, soulèvement du radier, les trois modes décrits par Wang et al.}
\label{fig-tun-tip16}
\end{figure}

\subsection{Amortissement et durée}
\label{sec-tunnel-duree}

Sur un maillage de mise au point et à $t = 0{,}10$~s, passer l'amortissement local de $0{,}15$ à $0{,}7$ divise par 4 le nombre de joints rompus ($3\,816$ contre 952), la longueur fissurée ($942$ contre 233~m) et le déplacement maximal ($0{,}069$ contre $0{,}017$~m), et réduit $R^{95}$ de $21~\%$. Sur une compression uniaxiale, le même changement relève le pic de $48~\%$. L'article ne chiffre pas son amortissement visqueux, de sorte que ce choix n'est pas arbitré par la référence. Le rejeu d'octobre confirme le facteur sur un maillage couvrant toute l'EDZ (\cref{sec-tunnel-rejeu}).

À $t = 0{,}25$~s, la paroi converge encore linéairement. Tous les calculs longs gagnent encore $15$ à $31~\%$ de convergence sur leur dernier cinquième, au-dessus du seuil de stabilisation de $5~\%$ (\cref{fig-tun-stab}). Le calcul de référence prolongé, arrêté à $0{,}687$~s, en gagne $17~\%$. Toute comparaison de motif se fait donc à temps égal depuis le relâchement, et les grandeurs publiées dépendent de la durée d'intégration, que l'article ne précise qu'en nombre de pas.

\begin{figure}[htbp]
\centering
\includegraphics[width=0.85\linewidth]{tunnel/fig_wall_stab}
\caption{Convergence de paroi par secteur, calcul de référence prolongé : la paroi converge encore à $0{,}69$~s ($17~\%$ du déplacement sur le dernier cinquième), d'où la comparaison à temps égal.}
\label{fig-tun-stab}
\end{figure}

\begin{figure}[htbp]
\centering
\includegraphics[width=\linewidth]{tunnel/fig_tunnel_synthese}
\caption{Synthèse du banc du tunnel : (a) rayon d'EDZ en fonction de $\sigma_0$, rockim et Wang et al. ; (b) part de propagation en fonction de $k_\text{tip}$ ; (c) joints rompus et blocs à $0{,}41$~s en fonction de $k_\text{tip}$ ; (d) rapports entre amortissements $0{,}7$ et $0{,}15$, maillage de mise au point.}
\label{fig-tun-synthese}
\end{figure}

\subsection{Rejeu du 6 octobre 2026 avec le binaire courant}
\label{sec-tunnel-rejeu}

Les chiffres précédents viennent de binaires d'août 2026, antérieurs aux correctifs des lois de décharge et de contact. Quatre variantes ont été rejouées avec le binaire courant (\cle{f0209ef}). Le maillage de production ne tenait pas dans le temps disponible ; un maillage réduit a été généré par l'outil du dépôt : $22\,730$ triangles, $h = 0{,}40$~m dans la zone fine. Ce maillage viole la règle d'objectivité ($h > \ell_{cz}/2 = 0{,}278$~m), et le générateur le signale. Le rejeu n'est donc qu'un contrôle de tendance. Les fichiers de calcul sont ceux de l'étude, avec seulement le maillage changé et la durée fixée à $0{,}23$~s après le début du relâchement pour toutes les variantes. L'amortissement fort est la référence avec un amortissement local de $0{,}7$.

\begin{table}[htbp]
\centering
\caption{Rejeu du 6 octobre 2026 (binaire \cle{f0209ef}, maillage réduit de $22\,730$ triangles), $0{,}23$~s après le début du relâchement. Blocs comptés dans un rayon de 25~m ; n.m. : non mesuré.}
\label{tab-tun-rejeu}
\footnotesize
\setlength{\tabcolsep}{4pt}
\begin{tabular}{@{}lcccc@{}}
\toprule
 & adaptative & $k_\text{tip} = 1{,}6$ & amortissement $0{,}7$ & intrinsèque \\
\midrule
joints rompus & $6\,397$ & $11\,846$ & $1\,663$ & $13\,674$ \\
traction / cisaillement & $2\,385 / 4\,012$ & $3\,553 / 8\,293$ & $596 / 1\,067$ & $879 / 12\,795$ \\
$R^{95}$ / $R^\text{max}$ (m) & $15{,}02 / 17{,}41$ & $19{,}71 / 27{,}66$ & $9{,}63 / 10{,}52$ & $24{,}17 / 54{,}15$ \\
demi-axes $a_h$ / $a_v$ (m) & $12{,}69 / 13{,}93$ & $16{,}28 / 17{,}01$ & $8{,}71 / 8{,}84$ & $20{,}28 / 18{,}75$ \\
longueur fissurée (m) & $2\,425$ & $4\,770$ & $637$ & $6\,979$ \\
paroi, moyenne / quantile $95~\%$ (m) & $0{,}094 / 0{,}128$ & $0{,}135 / 0{,}186$ & $0{,}025 / 0{,}035$ & $0{,}186 / 0{,}248$ \\
blocs / part d'un élément & $2\,483$ / $74{,}7~\%$ & $5\,765$ / $85{,}1~\%$ & n.m. & $6\,473$ / $79{,}5~\%$ \\
part de propagation & $29{,}2~\%$ & $42{,}3~\%$ & n.m. & $52{,}5~\%$ \\
temps de calcul (1 fil) & 7~min~38~s & 8~min~42~s & 4~min~34~s & 22~min~02~s \\
\midrule
travail net du contact (J/m) & $2{,}69\times10^4$ & $1{,}94\times10^6$ & $-2{,}51\times10^4$ & $1{,}89\times10^7$ \\
énergie cohésive (J/m) & $4{,}29\times10^5$ & $9{,}57\times10^5$ & n.m. & $2{,}66\times10^6$ \\
\bottomrule
\end{tabular}
\end{table}

Le rejeu confirme les tendances de la production (\cref{tab-tun-rejeu,fig-tun-rejeu,fig-tun-rejeu-histoire}).
\begin{enumerate}
\item La référence réduite redonne l'ordre de grandeur archivé : $R^\text{max} = 17{,}41$~m contre $17{,}19$~m en production et 19~m chez Wang, $R^{95} = 15{,}02$~m contre $14{,}49$~m, EDZ quasi circulaire. Le nombre de joints rompus n'est pas comparable, le maillage étant $4{,}7$ fois plus grossier.
\item Le facteur de pointe $1{,}6$ multiplie les joints rompus par $1{,}85$, étend l'EDZ ($R^{95}$ $+31~\%$) et fait passer la propagation de $29{,}2$ à $42{,}3~\%$ : même sens qu'en production, à des niveaux plus bas.
\item L'amortissement $0{,}7$ divise par $3{,}8$ les joints rompus, la longueur fissurée et la convergence de paroi : le facteur 4 du maillage de mise au point est confirmé.
\item L'intrinsèque rompt $2{,}1$ fois plus de joints que l'adaptative, à $94~\%$ en cisaillement, découpe des blocs de 40 à 120~m$^2$, et des fissures atteignent le bord du domaine ($R^\text{max} = 54$~m pour un demi-côté de 50~m). Le contraste de motif de la production est retrouvé, avec $52{,}5~\%$ de propagation contre $58{,}9~\%$.
\end{enumerate}

\begin{figure}[htbp]
\centering
\includegraphics[width=\linewidth]{tunnel/fig_tunnel_rejeu}
\caption{Rejeu sur maillage réduit (binaire \cle{f0209ef}), quatre variantes à $0{,}23$~s après le début du relâchement. En haut, norme du déplacement sur la configuration déformée et joints rompus (rouge : cisaillement ; blanc : traction). En bas, rein droit en rendu éléments, déplacements amplifiés 5 fois.}
\label{fig-tun-rejeu}
\end{figure}

\begin{figure}[htbp]
\centering
\includegraphics[width=0.75\linewidth]{tunnel/fig_tunnel_rejeu_histoire}
\caption{Nombre de joints rompus en fonction du temps depuis le début du relâchement (rejeu du 6 octobre 2026). La fissuration démarre aux deux tiers de la rampe et ne sature dans aucune variante à $0{,}23$~s.}
\label{fig-tun-rejeu-histoire}
\end{figure}

Le rejeu montre aussi que le travail net injecté par le contact croît avec la fragmentation. Sur le maillage réduit, il reste petit devant l'énergie cohésive en adaptatif ($2{,}7\times10^4$ contre $4{,}3\times10^5$~J/m). Il la dépasse avec le facteur de pointe ($1{,}94$ contre $0{,}96$~MJ/m) et surtout en intrinsèque ($18{,}9$ contre $2{,}66$~MJ/m, soit 7 fois). Ce point était ouvert depuis août : en production, le contact valait déjà les trois quarts de l'énergie de fissuration (416 contre 545~kJ/m). Tant qu'il n'est pas instruit, toutes les EDZ du banc sont conditionnelles, et celles des variantes très fragmentées, facteur de pointe et intrinsèque, sont suspectes d'être gonflées par une injection numérique. Quatre calculs de diagnostic, une clé changée à la fois, en ont cherché l'origine (\cref{sec-tunnel-diag}) : aucune clé ne l'explique seule. Le diagnostic 3 en identifie la cause : une erreur d'intégration du potentiel de contact, qui croît avec sa raideur normale et avec le pas de temps. Sur la coupe 2D, la décharge \cle{origin} sans cliquet a créé 115~J/m dans le poste des joints, création que le cliquet a supprimée (\cref{sec-coupe-rejeu}). Dans le tunnel, le cliquet ne réduit le travail net injecté par le contact que de $16{,}2~\%$ : cette décharge y contribue, mais n'en est pas la source principale. Le garde-fou d'énergie (\cle{budgetAbortPct}, \cref{sec-B4}) était désactivé dans les rejeux du tunnel, et l'énergie libérée par l'excavation, qui situerait ces travaux à l'échelle du problème, n'a pas été relevée. Le résidu du bilan d'énergie, entre $1{,}1$ et $41~\%$ selon la variante, dépasse le seuil de conformité de $1~\%$ (\cref{sec-B4}). Le diagnostic montre que ce résidu est divisé par 2 avec le pas de temps : c'est lui aussi une erreur d'intégration (\cref{sec-tunnel-diag}). Tant qu'elle n'est pas réduite, les travaux de contact du \cref{tab-tun-rejeu}, lus dans ce même bilan non fermé, ne sont pas vérifiés.

\subsection{Diagnostic de l'injection}
\label{sec-tunnel-diag}

Quatre calculs de diagnostic ont tourné dans la nuit du 6 au 7 octobre 2026 (binaire \cle{f638b0f}, 4 fils, environ 5~min chacun) sur la variante $k_\text{tip} = 1{,}6$, jusqu'à $t = 0{,}25$~s, soit $0{,}23$~s après le début du relâchement. Chacun ne change qu'une clé de cette base : naissance du contact en pénalité (\cle{gcBirth = penalty}), frottement de contact nul (\cle{contactMu = 0}), décharge \cle{plastic} au lieu d'\cle{origin}, et \cle{origin} avec cliquet (\cle{jointSecantRatchet = true}). L'amortissement local reste de $0{,}15$. Le maillage réduit de la section précédente n'étant pas archivé, il a été régénéré par l'outil du dépôt : $23\,004$ triangles, aussi grossier que l'original ($h > \ell_{cz}/2$). La base a été rejouée sur ce maillage. Elle garde la signature du rejeu : $2{,}61$~MJ/m de travail net du contact pour $0{,}998$~MJ/m d'énergie cohésive, soit $2{,}6$ fois (contre $1{,}94$ et $0{,}957$~MJ/m sur le maillage de $22\,730$ triangles). L'avertissement de non-conservation de la décharge \cle{origin} disparaît avec \cle{plastic} et avec le cliquet. Le tableau~\ref{tab-tun-diag} donne les bilans.

\begin{table}[htbp]
\centering
\caption{Diagnostic de l'injection par le contact, variante $k_\text{tip} = 1{,}6$, maillage régénéré de $23\,004$ triangles, $t = 0{,}25$~s. Une clé changée par calcul ; écart à la base entre parenthèses, sauf pour le résidu, rapporté à l'échelle du bilan ; vitesse maximale lue sur 14 trames.}
\label{tab-tun-diag}
\footnotesize
\setlength{\tabcolsep}{2.5pt}
\begin{tabular}{@{}p{0.21\linewidth}ccccc@{}}
\toprule
 & base & naissance & frottement & décharge & \cle{origin} \\
 & & en pénalité & nul & \cle{plastic} & avec cliquet \\
\midrule
contact, net injecté (MJ/m) & $2{,}611$ & $2{,}058$ ($-21{,}2~\%$) & $12{,}35$ ($+373~\%$) & $5{,}582$ ($+114~\%$) & $2{,}187$ ($-16{,}2~\%$) \\
joints, cohésif et amortisseur (MJ/m) & $1{,}075$ & $1{,}041$ & $1{,}508$ & $1{,}257$ & $1{,}011$ \\
cohésif (MJ/m) & $0{,}9983$ & $0{,}9679$ ($-3{,}0~\%$) & $1{,}379$ ($+38{,}2~\%$) & $1{,}178$ ($+18{,}0~\%$) & $0{,}9461$ ($-5{,}2~\%$) \\
frottement (MJ/m) & $2{,}708$ & $2{,}615$ & 0 & $4{,}059$ & $2{,}497$ \\
Cundall (MJ/m) & $9{,}127$ & $8{,}550$ ($-6{,}3~\%$) & $19{,}57$ ($+114~\%$) & $12{,}35$ ($+35{,}3~\%$) & $8{,}515$ ($-6{,}7~\%$) \\
résidu du bilan (MJ/m) & $-1{,}328$ ($6{,}62~\%$) & $-1{,}340$ ($7{,}15~\%$) & $-1{,}333$ ($3{,}37~\%$) & $-1{,}304$ ($5{,}02~\%$) & $-1{,}325$ ($7{,}02~\%$) \\
joints rompus & $11\,536$ & $11\,462$ & $12\,202$ & $11\,535$ & $11\,784$ \\
joints insérés & $13\,117$ & $13\,193$ & $13\,825$ & $13\,403$ & $13\,406$ \\
vitesse nodale max (m/s) & $25{,}85$ & $11{,}88$ & $23{,}80$ & $11{,}27$ & $11{,}18$ \\
\bottomrule
\end{tabular}
\end{table}

\begin{enumerate}
\item Aucune des quatre clés n'explique l'injection seule. La plus efficace, la naissance en pénalité, retire $21{,}2~\%$ du travail net injecté ; le contact injecte encore $2{,}1$ fois l'énergie cohésive.
\item Le frottement n'est pas la source, il la masque. Sans frottement, le travail net injecté par le contact croît de $373~\%$ ($12{,}35$~MJ/m) et l'amortissement de Cundall dissipe $114~\%$ de plus. Dans la base, le frottement dissipe $2{,}708$~MJ/m, autant que l'injection nette : la partie normale du contact injecterait donc de l'ordre de 5~MJ/m.
\item La décharge \cle{origin} n'est pas la source principale. Le cliquet ne réduit l'injection que de $16{,}2~\%$, et la décharge \cle{plastic} l'augmente de $114~\%$. Ce résultat réfute pour le tunnel l'hypothèse tirée de la coupe 2D, où le cliquet supprimait la création lue dans le poste des joints (\cref{sec-coupe-rejeu}) ; la contribution de cette décharge n'est pas nulle pour autant.
\item La naissance du contact sur un joint mort contribue : la naissance en pénalité retire $21{,}2~\%$ de l'injection et divise par deux la vitesse nodale maximale ($25{,}85$ puis $11{,}88$~m/s). La décharge \cle{plastic} et le cliquet divisent aussi cette vitesse par $2{,}3$ : le pic de la base est un événement local, pas l'injection globale.
\item Le résidu du bilan d'énergie varie peu entre ces cinq calculs, de $-1{,}30$ à $-1{,}34$~MJ/m (3 à $7~\%$), alors que les autres postes varient d'un facteur 2 à 5. Le diagnostic 3 montre qu'il n'est pas constant : il dépend du pas de temps.
\end{enumerate}

Un cinquième calcul combine frottement nul et naissance en pénalité (\cle{tip16_mu0_gcpen}, même maillage, 263~s), pour séparer la naissance du potentiel normal ; dans la base, $65~\%$ des joints meurent en compression. La naissance en pénalité retire la même part du travail net injecté avec ou sans frottement : $-21{,}2~\%$ avec, $-22{,}6~\%$ sans ($12{,}35$ puis $9{,}556$~MJ/m). Sans frottement et avec naissance en pénalité, le contact injecte encore $9{,}556$~MJ/m, soit $7{,}7$ fois l'énergie cohésive ($1{,}24$~MJ/m). La source dominante est donc la réponse normale du potentiel de contact elle-même, et non la naissance du contact ni la décharge des joints. La vitesse nodale maximale y atteint $30{,}1$~m/s et le résidu du bilan reste à $-1{,}328$~MJ/m. Le mécanisme diffère de celui de la coupe 2D, où la création se lisait dans le travail net des joints, et le lien entre les deux bancs reste une hypothèse (\cref{sec-coupe}).

Une troisième série de calculs, le diagnostic 3, a tourné le 7 octobre 2026 de 0~h~38 à 1~h~35, sur la base sans frottement (\cle{tip16_mu0}), pour isoler la partie normale du contact. Il change la raideur normale du potentiel $p$ (\cle{potPenaltyFactor} à $0{,}1$ et 10) ou divise le pas de temps par 2, sur le même maillage. Les trois calculs sont allés au bout ($t = 0{,}25$~s). Le tableau~\ref{tab-tun-diag3} donne les bilans.

\begin{table}[htbp]
\centering
\caption{Diagnostic 3, base sans frottement, même maillage, $t = 0{,}25$~s. Écart à la base entre parenthèses, sauf pour le résidu, rapporté à l'échelle du bilan ; vitesse maximale lue sur 14 trames.}
\label{tab-tun-diag3}
\footnotesize
\setlength{\tabcolsep}{3pt}
\begin{tabular}{@{}p{0.25\linewidth}cccc@{}}
\toprule
 & base & $p/10$ & $\Delta t/2$ & $10\,p$ \\
\midrule
$p$ (N/m) & $10^{10}$ & $10^{9}$ & $10^{10}$ & $10^{11}$ \\
$\Delta t$ (s) & $5{,}586\times10^{-6}$ & $5{,}586\times10^{-6}$ & $2{,}793\times10^{-6}$ & $3{,}484\times10^{-6}$ \\
contact, net injecté (MJ/m) & $12{,}35$ & $1{,}435$ ($-88{,}4~\%$) & $7{,}266$ ($-41{,}2~\%$) & $55{,}74$ ($+351~\%$) \\
cohésif (MJ/m) & $1{,}379$ & $0{,}8878$ ($-35{,}6~\%$) & $0{,}8630$ ($-37{,}4~\%$) & $2{,}186$ ($+58{,}5~\%$) \\
Cundall (MJ/m) & $19{,}57$ & $9{,}743$ ($-50{,}2~\%$) & $13{,}67$ ($-30{,}1~\%$) & $69{,}48$ ($+255~\%$) \\
résidu du bilan (MJ/m) & $-1{,}333$ ($3{,}37~\%$) & $-1{,}296$ ($5{,}98~\%$) & $-0{,}6543$ ($2{,}24~\%$) & $-0{,}8464$ ($0{,}66~\%$) \\
joints rompus & $12\,202$ & $11\,068$ & $11\,755$ & $13\,193$ \\
vitesse nodale max (m/s) & $23{,}80$ & $12{,}59$ & $22{,}83$ & $31{,}42$ \\
durée du calcul (s) & 311 & 690 & 1\,472 & 1\,272 \\
\bottomrule
\end{tabular}
\end{table}

\begin{enumerate}
\item L'injection suit la raideur normale du potentiel : $-88{,}4~\%$ à $p/10$, $+351~\%$ à $10\,p$, soit une pente log-log de $0{,}93$ puis $0{,}65$. Avec $10\,p$, le solveur réduit lui-même le pas de temps, ce qui n'empêche pas l'injection de quadrupler.
\item Elle suit aussi le pas de temps : $-41{,}2~\%$ à $\Delta t/2$, contre $-50~\%$ pour un biais d'ordre 1. C'est une erreur d'intégration du potentiel de contact, raideur mal résolue par le pas, et non un terme physique.
\item Le résidu du bilan n'est pas constant : il est divisé par 2 avec le pas de temps ($-1{,}333$ puis $-0{,}6543$~MJ/m). C'est lui aussi une erreur d'intégration, et non un travail de la précontrainte ou de l'excavation non compté.
\item À $p = 10^9$~N/m, le contact n'injecte plus que $1{,}6$ fois l'énergie cohésive ($1{,}435$ contre $0{,}8878$~MJ/m), contre 9 fois à $10^{10}$~N/m, pour un temps de calcul multiplié par $2{,}2$.
\item L'énergie cohésive et le nombre de joints rompus dépendent aussi de $p$ ($-36~\%$ et $-9~\%$ à $p/10$). Le faciès lui-même dépend de la raideur de contact : les faciès obtenus à des $p$ différents ne se comparent pas.
\end{enumerate}

Un diagnostic 4 vérifie $p/10$ avec frottement, sur la référence adaptative et sur la variante à facteur de pointe ; il a été lancé dans la nuit du 6 au 7 octobre et son résultat est à venir. La valeur $p = 10^9$~N/m n'est donc pas une recommandation validée, et l'interpénétration maximale à cette raideur n'a pas été relevée. Ces conclusions ont des réserves : maillage trop grossier, 4 fils (résultats non identiques au bit près), un seul tirage, durée de $0{,}25$~s seulement, vitesse maximale échantillonnée sur 14 trames et non à chaque pas.

\subsection{Volet exploratoire : tunnel dans un massif lité (méthode de Lisjak)}

Un second volet, exploratoire, reprend le même tunnel dans un massif lité à $45^\circ$. Il assemble les trois briques de Lisjak~\cite{lisjak2013thesis,lisjak2014shales} : élasticité transversalement isotrope, loi cohésive directionnelle et maillage à arêtes alignées sur le litage. Les paramètres dans le plan du litage sont ceux de Wang. Les rapports d'anisotropie sont ceux de l'argilite d'Opalinus : $E' = 0{,}342\,E$, $\nu' = 0{,}25$, $G' = 2{,}37$~GPa, et rapports cohésifs $0{,}246$ ($f_t$), $0{,}111$ ($c$), $0{,}057$ ($G_I$), $0{,}286$ ($G_{II}$). Chaque propriété cohésive varie linéairement avec l'angle $\gamma$ entre le joint et le litage :
\begin{equation}
X(\gamma) = X_\text{min} + (X_\text{max} - X_\text{min})\,\frac{\gamma}{90^\circ}, \qquad X \in \{f_t, c, G_I, G_{II}\} .
\end{equation}
Les contrôles élémentaires sont satisfaits : modules apparents à $0{,}13~\%$ près, rapport d'anisotropie du pic de $1{,}94$ (Lisjak : environ $1{,}9$).

Le seul calcul de production (maillage lité de $117\,132$ triangles, 4~h~11) donne une EDZ en losange orientée par le litage. Comparé au témoin isotrope à champ corrélé, il rompt $31~\%$ de joints en moins. La profondeur endommagée vaut $6{,}0$~m le long du litage et $3{,}4$~m en travers, soit un rapport de $1{,}77$ contre $1{,}19$ en isotrope. Aucune fissure ne traverse un plan délaminé, ce que prédit le critère de déviation de He et Hutchinson pour un rapport de ténacités interface sur roche de $0{,}057$, quatre fois sous le seuil d'environ $1/4$. Les réserves sont lourdes : un seul calcul, une graine et un pendage ; zone cohésive non résolue le long du litage ($\ell_{cz} = 179$~mm pour des arêtes de 200~mm) ; convergence de $3{,}4$~m aux reins, soit un effondrement. Aucune figure de résultat n'est archivée dans les dépôts. Les calculs aux pendages $0^\circ$ et $90^\circ$ sont préparés et n'ont pas été lancés.

\subsection{Synthèse}

Avec les paramètres cohésifs publiés, rockim reproduit la physique globale du banc (\cref{fig-tun-synthese}), avec un amortissement local de $0{,}15$ choisi par rockim qui pilote le nombre de joints rompus d'un facteur 4. Il retrouve un rayon d'EDZ maximal à $-4$ à $-21~\%$ de l'article entre 4 et 7~MPa et à $+15~\%$ à 3~MPa ; au quantile $95~\%$, les écarts valent $-19$ à $-31~\%$ et $-1~\%$. Il retrouve aussi la prédominance du cisaillement, des fissures conjuguées près de l'angle de Mohr-Coulomb, le basculement de l'ellipse avec $\lambda$ et le rapport $2{,}7$ sur le nombre de fissures. Le maximum est l'estimateur le plus favorable à l'accord. La précontrainte et l'excavation sont vérifiées par rapport aux solutions analytiques à $2{,}1~\%$ près. Le désaccord porte sur la texture : l'insertion adaptative granule un anneau là où MultiFracS découpe des blocs. Le schéma d'insertion en est la cause la plus probable, et le désordre des résistances n'y change rien. Cette conclusion reste conditionnelle : les variantes qui découpent des blocs, intrinsèque et facteur de pointe, sont aussi celles où le contact injecte le plus d'énergie, jusqu'à 7 fois l'énergie cohésive en intrinsèque ($18{,}9$ contre $2{,}66$~MJ/m). Le diagnostic de cette injection en identifie la cause : une erreur d'intégration du potentiel de contact, qui suit sa raideur normale ($-88{,}4~\%$ à $p/10$, $+351~\%$ à $10\,p$) et le pas de temps ($-41{,}2~\%$ à $\Delta t/2$). La décharge \cle{origin} n'en est pas la cause principale, le frottement la masque et la naissance du contact sur un joint mort y contribue (\cref{sec-tunnel-diag}). L'énergie cohésive et le nombre de joints rompus dépendent aussi de la raideur, et la raideur $p = 10^9$~N/m attend sa vérification avec frottement. Le rayon d'EDZ lui-même provient de calculs à décharge non conservative, à bilan d'énergie non fermé et à travail de contact comparable à l'énergie de fissuration : l'accord avec Wang est conditionnel au même diagnostic.

Les grandeurs scalaires publiées dépendent de choix que l'article ne documente pas : amortissement (facteur 4), durée d'intégration (la paroi converge encore), schéma d'insertion (facteur $1{,}7$ sur les joints rompus). La comparaison est donc qualitative et semi-quantitative ; ce n'est pas une validation. Restent ouverts la correction de l'injection d'énergie par le contact, artefact de raideur aggravé dans les calculs très fragmentés, la réduction de l'erreur d'intégration qui empêche le bilan d'énergie de fermer, la saturation de l'EDZ au-delà de 5~MPa, et le rejeu en taille réelle avec le binaire courant et une décharge conservative.
